% ECONOMICS_PROBLEM: {"id": "Q54-602", "title": "Barriers to adaptation", "jel_code": "Q54", "source_id": "OP-0434"}
% !TeX program = xelatex
\documentclass[11pt,a4paper]{article}
\usepackage[margin=23mm]{geometry}
\usepackage{fontspec}
\setmainfont{Latin Modern Roman}
\usepackage{amsmath,amssymb,mathtools}
\usepackage{xcolor,enumitem,booktabs,longtable,array,xurl,hyperref}
\definecolor{muted}{HTML}{53636C}
\hypersetup{colorlinks=true,urlcolor=blue,linkcolor=blue}
\setlength{\parindent}{0pt}
\setlength{\parskip}{5pt}
\newcommand{\R}{\mathbb R}
\newcommand{\E}{\mathbb E}
\newcommand{\N}{\mathbb N}
\newcommand{\ind}{\mathbf 1}
\newcommand{\Var}{\operatorname{Var}}
\newcommand{\Cov}{\operatorname{Cov}}
\newcommand{\supp}{\operatorname{supp}}
\DeclareMathOperator*{\argmin}{arg\,min}
\DeclareMathOperator*{\argmax}{arg\,max}
\newcommand{\entrymeta}[3]{{\sffamily\small\color{muted}\textbf{\texttt{#1}}\quad|\quad #2\par #3\par}}
\newcommand{\Problemref}[1]{\texttt{#1}}
\newcommand{\JELref}[2]{\texttt{#2} (JEL \texttt{#1})}
\begin{document}
\section*{Barriers to adaptation}
\textbf{Problem ID:} \texttt{Q54-602}\quad \textbf{JEL:} \texttt{Q54}\par
% BEGIN ECONOMICS STATEMENT
\label{problem:OP-0434}
{\sffamily\small\color{muted}Original subject group: Climate, environment and energy\par}
\label{definitions:problem:30007716}
\textbf{Audit classification:} Published research agenda. Checked source identity and appropriate agenda/background support; expanded intervention, units, population, definedness and causal restrictions. Exact targets and variants are editorial.
\subsection*{Definitions and precise research target}
\textit{Editorial formalization.} Consider farmers in drought-prone districts of one low-income country offered a specified water-saving technology. Let \(I_i,C_i,R_i\in\{0,1\}\) indicate randomized offers of verified information, affordable credit, and drought insurance to farmer \(i\). Let \(D_i(I_i,C_i,R_i)\) denote adoption within two seasons, and \(V_i(I_i,C_i,R_i)\) denote five-year household welfare net of technology costs. Welfare combines a stated consumption utility measure with agricultural risk exposure. For every offer combination \(z\in\{0,1\}^3\), define
\[\tau_D(z)=E[D_i(z)-D_i(0,0,0)],\qquad\tau_V(z)=E[V_i(z)-V_i(0,0,0)].\]
Expectations concern eligible baseline farmers, including nonadopters. Identify these contrasts and the interactions among offers under consistency, randomized assignment, and either negligible cross-farm spillovers or an explicitly modeled district exposure measure. Then determine which adoption barriers are policy-relevant at the specified prices and how their welfare effects change under a stated future drought distribution. Offer effects do not uniquely reveal psychological or financial mechanisms without additional exclusions. Distinguish a package that increases adoption from one that increases welfare, and measure crowd-out of existing insurance or investment.

Assignment \(Z_i\in\{0,1\}^3\) has positive probability in every arm and is independent of all potential outcomes, conditional on declared randomization strata. Adoption \(D_i(z)\in\{0,1\}\) is use of the specified technology by the second season, rather than accepting an offer. Define \(V_i(z)=E[\sum_{t=0}^{4}\beta^tu(c_{it}(z))\mid i]\), with technology costs included in the consumption budget; risk is thereby valued through expected utility. Require integrable welfare. A pairwise interaction, for example, is \(\Gamma_{IC}(r)=E[V_i(1,1,r)-V_i(1,0,r)-V_i(0,1,r)+V_i(0,0,r)]\). These offer interactions do not identify separate latent constraints. Under interference replace \(V_i(z)\) by \(V_i(z,g)\), with district exposure \(g\) and its randomization explicitly specified. All expectations use a stated baseline population and probability law. Identification means every admissible data-generating model with the same observed-data law yields the same displayed target; otherwise report the identified set. Exact equations, horizons and assumptions are editorial, rather than source-given canonical conjectures.
\subsection*{Variants and assumption changes}
\begin{enumerate}
\item \textbf{Editorial adoption variant.} Use randomized offers as instruments for actual use only under exclusion, monotonicity and a positive first stage. The complier effect differs from all-eligible offer effects.
\item \textbf{Editorial spillover variant.} Randomize district saturation \(s\) and personal offers \(z\). Identify \(E[V(z,s)-V(0,s)]\) and \(E[V(0,s)-V(0,0)]\) under between-district noninterference.
\end{enumerate}
\subsection*{References and source audit}
\textbf{1.} Explicit policy-oriented adaptation research agenda; exact interventions and optimizers are editorial. Locator: VoxEU column, 15 April 2025; introductory summary. Primary indexed author summary checked; subsequent full-text retrieval failed, so precise subsection claims are not independently reaffirmed. URL: \url{https://cepr.org/voxeu/columns/economics-climate-adaptation-academic-insights-effective-policy}.
\begin{enumerate}\raggedright
\item\hypertarget{definitions:source:30007716:1}{} Tamma Carleton; Esther Duflo; Kelsey Jack; Guglielmo Zappalà. 2025. \emph{The economics of climate adaptation: From academic insights to effective policy}. VoxEU column, 15 April 2025; introductory summary.
\url{https://cepr.org/voxeu/columns/economics-climate-adaptation-academic-insights-effective-policy}.
\end{enumerate}

\subsection*{Common conventions: Source mathematical conventions}

These conventions apply to each entry unless explicitly replaced.

\textbf{Models and identification.} A primitive model \(m\) specifies all probability laws, feasible choices, information, equilibrium and measurement rules used by its target. The admissible class \(\mathcal M\) is fixed independently of outcomes. If \(P_O(m)\) is its observable probability law and \(\tau(m)\) the target, its identified set at observable law \(P\) is
\[
 \mathcal I_\tau(P)=\{\tau(m):m\in\mathcal M,\ P_O(m)=P\}.
\]
Point identification means that this set is a singleton. An empty set signals incompatibility of data and assumptions. Coordinate bounds need not describe the joint set. Bounds are sharp only if every retained target is attainable or arbitrarily approachable by compatible models. Finite bounds require explicit restrictions on unobserved quantities; unrestricted confounding, hidden wealth, extrapolation or arbitrary equilibrium selection can yield uninformative sets.

\textbf{Causal interventions.} A potential outcome \(Y_i(p)\) is the outcome under a complete policy regime \(p\), including timing, financing, information and market spillovers. The target population and units are fixed before treatment. With interference, use the complete assignment vector or a justified exposure mapping; individual no-interference notation alone cannot identify an economy-wide intervention. A difference of mean potential outcomes does not require the cross-world joint distribution. Conditioning on post-treatment migration, survival, enrollment or employment changes the estimand and needs separate justification.

\textbf{Statistical statements.} An estimator, confidence set, forecast or test specifies a sampling law, model class and loss/coverage criterion. Uniform \((1-\alpha)\)-coverage means \(\inf_{m\in\mathcal M}\Pr_m\{\tau(m)\in C_n\}\geq1-\alpha\), or an explicitly asymptotic version. The whole parameter space trivially has coverage, so informativeness requires a stated width, risk or power objective. A forecasting improvement is relative to a named benchmark, information set and evaluation population.

\textbf{Equilibrium and optimization.} An equilibrium comprises feasible plans, optimality, beliefs where needed, budget constraints and all market/resource clearing. The primitives or a stated benchmark define each correspondence \(\mathcal E(p;m)\). Nonemptiness is assumed or proved, never inferred just from compact policy sets. A selected-equilibrium target uses a declared selection rule; a robust target quantifies over all permitted equilibria. Use suprema or infima unless attainment is established. An \(\varepsilon\)-optimal policy is feasible and within \(\varepsilon\) of the finite optimum value. Report an empty feasible set separately.

\textbf{Welfare and accounting.} Normative utility, interpersonal normalization, social weights, numeraire, discounting and the use of public receipts are primitives. Behavioral utility need not coincide with the welfare criterion. Household welfare includes ownership income and rents once; adding profits again to compensating variation generally double-counts them. A monetary equivalent is defined by an explicit transfer and price convention, with monotonicity and range conditions. Average effects, mechanism contributions, transfers and welfare losses are different targets. Infinite sums require convergence and dynamic feasible plans require terminal or no-Ponzi conditions.

\textbf{Source versions.} A working-paper date, revision date and journal publication year are distinguished. A DOI for a journal version is not automatically the DOI of an earlier manuscript. Locators refer to the stated version and printed pages where specified. A metadata-only check does not verify a claim about a full-text conclusion. Every entry's source subsection narrows the provenance claim accordingly.
% END ECONOMICS STATEMENT
\end{document}
